\section{A15.1b2d core-both closure and full Sub-wedge A}
\label{sec:p12-b2d-core-both}

We close the only remaining part of b2d.  Throughout this section let
\[
e/2 \le R < d/2,\qquad
R < \sigma < \varepsilon < \varepsilon_{\max},
\qquad
\sigma>d/2,
\]
and let \(h\in L^2(R,S)\) satisfy \(L_{T_0}h=0\), with
\(S=T+\sigma\), \(T_0=T+\varepsilon\).
Set
\[
I_{\rm both}:=
\bigl(\max\{R,d-\sigma\},\,\min\{\sigma,d-R\}\bigr).
\]
For \(x\in I_{\rm both}\) write
\[
y:=d-x,\qquad \kappa:=e-\delta.
\]
Then \(x,y\in(R,\sigma)\), and \(I_{\rm both}\) is invariant under
\(x\mapsto d-x\).

\begin{lemma}[Uniform core-both support bounds]
\label{lem:p12-b2d-core-both-support}
For a.e.\ \(x\in I_{\rm both}\),
\[
0<e-x<R,\qquad 0<e-y<R,
\]
\[
0<x-\delta<R,\qquad 0<y-\delta<R,
\]
and
\[
0<2\delta-x<R.
\]
Moreover
\[
y+\delta>\sigma,\qquad x+\kappa>\sigma.
\]
\end{lemma}

\begin{proof}
Since \(x,y\in(R,\sigma)\), \(R\ge e/2\), and
\(\sigma<\varepsilon_{\max}<e\),
\[
0<e-x<e-R\le e/2\le R,
\]
and the same holds with \(y\) in place of \(x\).

Next \(x,y>R\ge e/2>\delta\), while
\[
x-\delta,\ y-\delta
<\varepsilon_{\max}-\delta<e/2\le R.
\]
Also \(\varepsilon_{\max}<2\delta\), so \(2\delta-x>0\), and
\(R>\delta\) gives
\[
2\delta-x<2\delta-R<R.
\]

Finally
\[
y+\delta>e/2+\delta>\varepsilon_{\max}>\sigma.
\]
Since \(e>2\delta\), \(\kappa=e-\delta>\delta\), hence also
\[
x+\kappa>e/2+\delta>\varepsilon_{\max}>\sigma.
\]
The elementary comparisons used here are
\(e>2\delta\iff256>243\),
\(\varepsilon_{\max}<e\iff15<16\),
\(\varepsilon_{\max}<2\delta\iff80<81\), and
\(\varepsilon_{\max}-\delta<e/2\iff100<108\).
\end{proof}

\begin{theorem}[Full b2d core-both local kill]
\label{thm:p12-b2d-core-both-full}
Under the hypotheses above,
\[
h(x)=H(x)=l(x)=0
\]
for a.e.\ \(x\in I_{\rm both}\).
\end{theorem}

\begin{proof}
Each source map below is affine in \(x\) with derivative \(\pm1\).
Hence the preimage of a source-null set is null, and after intersecting
finitely many full-measure sets all source equations may be used
simultaneously for a.e.\ \(x\in I_{\rm both}\).

Consider the thirteen source positions
\[
\begin{aligned}
u_1&=d-x, &
u_2&=2d-x, &
u_3&=b-x, &
u_4&=3d-x,\\
u_5&=T-x, &
u_6&=T+\delta-x, &
u_7&=T+d-x,\\
u_8&=e+x, &
u_9&=a+x, &
u_{10}&=3e+x,\\
u_{11}&=b-\delta+x, &
u_{12}&=T-\delta+x, &
u_{13}&=T+x.
\end{aligned}
\]
They are all horizon-legal.  Indeed, the only sources above \(T\) are
\[
u_7=T+y<T+\sigma<T_0,\qquad
u_{12}=T+(x-\delta)<T+x<T+\sigma<T_0,
\]
and
\[
u_{13}=T+x<T+\sigma<T_0.
\]
All remaining sources are \(<T\): for example
\(u_2=d+y<d+e=a\),
\(u_4=2d+y<2d+e=b\),
\(u_6<T\) because \(x>\delta\),
\(u_9<a+e<T\), and
\(u_{10}<3e+\varepsilon_{\max}<4e<T\).

Applying
\[
Lh(u)=p[h(u-a)-h(u+a)]
+r[h(u-b)-h(u+b)]
+q[h(u-T)-h(u+T)]
\]
with the anti-reflection \(h(-t)=-h(t)\), zero extension outside
\((R,S)\), and Lemma~\ref{lem:p12-b2d-core-both-support}, gives:
\begin{align*}
Q_1:\;&-p h(e+x)-p h(b-x)-r h(a+x)-r h(T+\delta-x)\\
&\hspace{28mm}-q h(b-\delta+x)-q H(y)=0,\\
Q_2:\;&-p h(T+\delta-x)-r h(e+x)-q h(2e+x)=0,\\
Q_3:\;&p h(y)-pH(y)-r h(x)-q h(e+x)=0,\\
Q_4:\;&p h(y+\delta)-q h(x+\kappa)=0,\\
Q_5:\;&p h(a-x)-q h(x)=0,\\
Q_6:\;&p h(2d-x)+r h(y)=0,\\
Q_7:\;&p h(b-x)+r h(a-x)+q h(y)=0,\\
Q_8:\;&-p h(y)-p l(y)-r h(2d-x)-rH(x)-q h(b-x)=0,\\
Q_9:\;&p h(x)-pH(x)-r h(y)-q h(a-x)=0,\\
Q_{10}:\;&p h(x+\kappa)-q h(y+\delta)=0,\\
Q_{11}:\;&p h(e+x)-q h(y)=0,\\
Q_{12}:\;&p h(2e+x)+r h(x+\kappa)=0,\\
Q_{13}:\;&p h(a+x)+r h(e+x)+q h(x)=0.
\end{align*}
Here \(h(b-\delta+x)=h(T-y)=l(y)\).

Let \(\Delta:=p^2-q^2\).  From \(Q_4,Q_{10}\),
\[
\Delta\,h(y+\delta)=\Delta\,h(x+\kappa)=0.
\]
Since \(\Delta>0\),
\[
h(y+\delta)=h(x+\kappa)=0.
\]
Then \(Q_{12}\) gives \(h(2e+x)=0\).  The remaining equations yield
successively
\[
h(e+x)=\frac qp\,h(y),
\qquad
h(T+\delta-x)=-\frac{rq}{p^2}h(y),
\]
\[
h(2d-x)=-\frac rp\,h(y),
\qquad
h(a-x)=\frac qp\,h(x),
\]
\[
h(a+x)=-\frac qp\,h(x)-\frac{rq}{p^2}h(y),
\]
\[
h(b-x)=-\frac{rq}{p^2}h(x)-\frac qp\,h(y),
\]
\[
H(y)=\frac{\Delta}{p^2}h(y)-\frac rp\,h(x),
\qquad
H(x)=\frac{\Delta}{p^2}h(x)-\frac rp\,h(y).
\]
Substitution in \(Q_8\) gives
\[
l(y)=
-\frac{r(p^2-2q^2)}{p^3}h(x)
-\frac{\Delta-2r^2}{p^2}h(y).
\]
Finally \(Q_1\) reduces exactly to
\[
\frac{2qr(2p^2-q^2)}{p^3}\,h(x)=0.
\]
Since \(p,q,r>0\) and \(2p^2-q^2>0\),
\[
h(x)=0
\]
for a.e.\ \(x\in I_{\rm both}\).

The involution \(x\mapsto y=d-x\) preserves \(I_{\rm both}\) and
Lebesgue-null sets, so the same conclusion gives \(h(y)=0\).
Then \(Q_5,Q_9\) give \(h(a-x)=H(x)=0\), and the same argument at
\(y\) gives \(H(y)=0\).  Identity P1
(Theorem~\ref{thm:p12-P1}) therefore gives \(l(x)=0\).
\end{proof}

\begin{remark}[Polynomial row certificate]
\label{rem:p12-b2d-core-both-certificate}
The thirteen rows admit an independent exact certificate.  With
\(\Delta=p^2-q^2\), set
\[
\begin{aligned}
\widetilde C=(&
p^3\Delta,\,
-p^2r\Delta,\,
-p^2q\Delta,\,
pq^2r^2,\,
-r\Delta(p^2-2q^2),\\
&-pqr\Delta,\,
p\Delta^2,\,
-p^2q\Delta,\,
pqr\Delta,\,
p^2qr^2,\\
&p\Delta(\Delta-2r^2),\,
-pqr\Delta,\,
p^2r\Delta).
\end{aligned}
\]
Then, identically in all visibility variables,
\[
\sum_{j=1}^{13}\widetilde C_jQ_j
=
2qr\Delta(2p^2-q^2)\,h(x).
\]
Thus the triangular elimination and the row-multiplier calculation
cross-check the same non-vanishing closure.
\end{remark}

\begin{theorem}[Full A15.1b2d Sub-wedge A injectivity]
\label{thm:p12-b2d-full}
Let
\[
2a<T_0<c,\qquad e/2\le R<d/2,\qquad T<S<T_0.
\]
Then
\[
\ker L_{R,S,T_0}^{\{a,b,2a\}}=\{0\}.
\]
\end{theorem}

\begin{proof}
Write \(\sigma=S-T\).

If \(\sigma\le d/2\), this is
Corollary~\ref{cor:p12-b2d-wedge-plus}.  Assume therefore
\(\sigma>d/2\).

If \(d/2<\sigma\le d-R\), the defect strip decomposes, up to endpoints,
as
\[
(R,d-\sigma)\quad\text{(core-single)}
\]
and
\[
(d-\sigma,\sigma)\quad\text{(core-both)}.
\]
The first interval is killed by
Theorem~\ref{thm:p12-b2d-core-single-full}, including its visibility
values \(H,l\); the second is killed by
Theorem~\ref{thm:p12-b2d-core-both-full}.

If \(\sigma>d-R\), the defect strip decomposes as
\[
(R,d-R)\quad\text{(core-both)}
\]
and
\[
(d-R,\sigma)\quad\text{(b2d-upper)}.
\]
The first interval is killed by
Theorem~\ref{thm:p12-b2d-core-both-full}; the second by
Corollary~\ref{cor:p12-b2d-upper-local}.  Hence in either case
\[
h=H=l=0
\quad\text{a.e. on }(R,\sigma).
\tag{C}
\]

Consequently the \(E_\sigma\)-equation is homogeneous on all of
\((R,a)\).  The reflection Steps~1--5 of
Theorem~\ref{thm:p12-b2b} apply exactly as in the proof of
Corollary~\ref{cor:p12-b2d-wedge-plus}, and give
\[
h=0
\quad\text{a.e. on }(R,a).
\tag{E}
\]

It remains to kill the small tail.  Let \(0<t<R\).
In P2,
\[
D(t)=(p^2-q^2-r^2)h(t)+qr[h(e-t)-h(e+t)].
\]
Here \(h(t)=0\) by support, \(e+t\in(R,a)\) and is zero by (E), while
\(e-t\) is either below \(R\) or belongs to \((R,e)\subset(R,a)\).
Thus \(D(t)=0\), so P2 gives
\[
H(t)=l(t).
\tag{1}
\]

For P1 put \(z=d-t\).  Since
\[
z>d-R>R,
\]
either \(z\ge\sigma\), in which case \(H(z)=0\) by tail support, or
\(R<z<\sigma\), in which case \(H(z)=0\) by (C).  Hence always
\(H(d-t)=0\).  P1 gives
\[
H(t)+l(t)=0.
\tag{2}
\]
From (1)--(2),
\[
H(t)=l(t)=0
\quad\text{for a.e. }0<t<R.
\]
Together with (C), the entire tail \((T,S)\) vanishes.  The support
therefore reduces to \((R,T)\), and
Theorem~\ref{thm:p12-b1} at \((R,T,T_0)\) gives \(h=0\).
\end{proof}
